% =============================================================================
\section{Prompts}
\label{app:prompts}
% =============================================================================

This appendix contains the system prompts used in the pipeline
experiments. Each prompt consists of a role/goal statement,
step-specific instructions, and (where applicable) the shared formalism
context. The formalism context is identical across the entity
extraction, constraint formalization, and chain-of-thought prompts; it
is shown once below and omitted from subsequent prompts for brevity.

% -----------------------------------------------
\subsection{Shared Formalism Context}
\label{app:formalism-context}

The following is an abbreviated version of the formalism context block
included in the entity extraction, constraint formalization, and CoT
prompts. The actual prompt lists each entity type with its permitted
operations, each temporal construct with its parameter names, each
predicate with its field structure, and each expression type with its
arguments, along with a worked example of a complete formalization.

{\footnotesize\ttfamily
\begin{quote}
ENTITY TYPES:
Signal (calculate), Storage (write, read),
EventTrigger (fire), Event (linked via modifiers),
Channel (transmit, receive), State (set\_state),
Value, Abstract, Set (insert, remove, clear)

TEMPORAL CONSTRUCTS:
Always, Eventually, Initial, Causes, CausesWithin,
Sequence, Precedes, Excludes, Immediately

PREDICATES:
Cmp, AllOf, AnyOf, Implies, Not, ForAll, Exists,
Happening, HasHappened, Operation

EXPRESSIONS:
Constant, ArithOp, Now, Prev, Next,
Val, ValBefore, EvtOccCount, LastOcc, FirstOcc,
MaxVal, MinVal, MkInterval, Start, End, Duration,
Diff, Addr, Size, Filter

IMPORTANT PATTERNS:
EventTrigger + Event pairs: when an operation fires
on an entity, corresponding Event entities (linked
via target and type modifiers) become true.
Time is ambient: expressions use Now() and temporal
constructs supply the time context.
\end{quote}
}

Additionally, the prompts include a worked example
showing the formalization of a scheduling deadline requirement
(CausesWithin with Happening predicates).

% -----------------------------------------------
\subsection{Ambiguity Review}
\label{app:prompt-ambiguity}

{\footnotesize\ttfamily
\begin{quote}
You are Ambiguity Reviewer. You are an expert in
analyzing safety-critical avionics requirements for
ambiguity and potential misinterpretation.

Review requirements against the ambiguity checklist
and flag genuine issues without rewriting the
requirement.

Checklist: dangling else / scope of action, ambiguity
of reference, omissions (causes without effects,
effects without causes), ambiguous logical operators,
negation scope, vague verbs/adverbs/adjectives,
built-in assumptions, ambiguous precedence, implicit
cases, temporal ambiguity, boundary ambiguity.

Do NOT rewrite or modify the requirement.
\end{quote}
}

% -----------------------------------------------
\subsection{Flavour Extraction}
\label{app:prompt-flavour}

{\footnotesize\ttfamily
\begin{quote}
You are a Time Model Classifier.

Discrete: step-by-step behavior, register operations,
state transitions, sequential actions without
real-time durations.

Continuous: real-time durations (seconds, ms, ns),
deadlines, continuous physical quantities.
\end{quote}
}

% -----------------------------------------------
\subsection{Entity Extraction}
\label{app:prompt-entity}

Uses the shared formalism context
(Appendix~\ref{app:formalism-context}), followed by:

{\footnotesize\ttfamily
\begin{quote}
You are Entity Extractor.


Each entity has: name (unique identifier), type,
modifiers (key-value pairs), description.
Every entity name must be unique.

IMPORTANT: For every EventTrigger, create a
corresponding Event entity with target and type
modifiers. For operations that produce events, create
corresponding Event entities if the requirement
references those occurrences.
\end{quote}
}

% -----------------------------------------------
\subsection{Constraint Formalization (Pipeline)}
\label{app:prompt-formula}

Uses the shared formalism context, followed by a
step-by-step decomposition algorithm:

{\footnotesize\ttfamily
\begin{quote}
You are Requirement Formalizer.


STEP 1 -- IDENTIFY TOP-LEVEL TEMPORAL SCOPE:
Always, Eventually, Causes/CausesWithin, Sequence,
Initial, or TraceAllOf. Separate condition from
effect.

STEP 2 -- DECOMPOSE CONDITION:
Split into sub-predicates. Event happening?
Value comparison? Occurrence count? Past event
reference using LastOcc/Start?

STEP 3 -- DECOMPOSE EFFECT:
What value must the entity have? Toggle (use
ValBefore)? Just an event firing (use Happening)?

STEP 4 -- NORMALIZE:
Entity names must match exactly. Operations must be
compatible with entity types. Time references use
Now(), Prev(Now()), Start(LastOcc(...)), etc.
\end{quote}
}

% -----------------------------------------------
\subsection{Chain-of-Thought (Single Call)}
\label{app:prompt-cot}

The CoT variant combines flavour extraction, entity extraction, and
constraint formalization into a single prompt. It uses the same
formalism context and adds a two-stage decomposition:

{\footnotesize\ttfamily
\begin{quote}
STAGE 1: MACRO-STRUCTURE

Step 1: Determine time model (Discrete/Continuous).
Step 2: Identify top-level temporal scope and
separate condition from effect.

STAGE 2: RECURSIVE DECOMPOSITION

Step 3: Extract all entities with types and modifiers.
Step 4: Decompose condition into sub-predicates.
Step 5: Decompose effect.
Step 6: Normalize entity references and operations.

Output the complete Requirement as JSON.
\end{quote}
}
